\chapter{Sobolev compactness on closed manifolds}\label{sec:sobolev-closed-compactness}
The mollification argument used for BV compactness in Theorem \ref{thm:bv-local-compactness} also gives the $L^2$ compactness needed for fixed-parameter elliptic problems. The translation-and-smoothing method parallels \cite{Sim18}, Chapter 2, Theorem 2.6; the phase-field application is the compactness step in \cite{ChAC}, Exercise 4.4. We give the $H^1$ version explicitly.

\begin{lemma}[The quadratic translation estimate]\label{lem:sobolev-l2-translations}
For $f\in H^1(\mathbb R^n)$ and $h\in\mathbb R^n$,
\[
 \|f(\,\cdot+h)-f\|_{L^2}\leq |h|\,\|Df\|_{L^2}.
\]
If $\rho_\delta$ is a nonnegative smooth mollifier of integral one supported in $B_\delta$, then
$\|\rho_\delta*f-f\|_2\leq\delta\|Df\|_2$.
\end{lemma}
\begin{proof}
For smooth $f$, the fundamental theorem of calculus gives
$f(x+h)-f(x)=\int_0^1 Df(x+sh)\cdot h\,ds$.
Square, apply Cauchy--Schwarz in $s$, and integrate in $x$; translation invariance of Lebesgue measure gives the first estimate. For $H^1$ functions apply the calculation to their smooth mollifications and pass to the limit in $L^2$ for the functions and their weak derivatives. Approximation by mollifiers follows from translation continuity in $L^2$; it does not require any compactness theorem. Finally average the first estimate against $\rho_\delta(h)$ and use the triangle inequality for the $L^2$ integral to obtain the second.
\end{proof}

\begin{theorem}[Strong L2 compactness on a closed manifold]\label{thm:sobolev-closed-l2-compactness}
On a smooth closed Riemannian manifold $(M^n,g)$, every bounded sequence in $H^1(M)$ has a strongly $L^2(M)$ convergent subsequence. The same is true for an $H^1(\mathbb R^n)$ bounded sequence whose members are supported in one fixed compact set. This is the compact embedding used in the fixed-parameter problem of \cite{ChAC}, Exercise 4.4.
\end{theorem}
\begin{proof}
First take the Euclidean sequence with common compact support $K$. For fixed $\delta>0$, its convolutions have support in $K+\overline B_\delta$ and uniform bounds on values and first derivatives:
\[
 \|\rho_\delta*f_j\|_\infty\leq\|\rho_\delta\|_2\|f_j\|_2,
 \qquad
 \|D\rho_\delta*f_j\|_\infty\leq\|D\rho_\delta\|_2\|f_j\|_2.
\]
Arzela--Ascoli makes this family precompact uniformly on the common compact set, hence in $L^2$. Lemma \ref{lem:sobolev-l2-translations} bounds the $L^2$ distance from each original function to its convolution by $C\delta$. For each tolerance choose $\delta$ small and then a finite net for the convolutions. The same net, with the uniformly small smoothing error added, is a finite net for the original sequence. Total boundedness and completeness of $L^2$ yield a convergent subsequence.

On $M$, choose finitely many coordinate charts and a smooth partition of unity with each cutoff supported compactly in its chart. The coordinate representations of the cutoff functions times $f_j$, extended by zero, have a common compact support and uniform Euclidean $H^1$ bounds. This follows from the product rule and bounded comparability of the coordinate metric and volume density on the compact supports; zero extension creates no boundary jump there. Apply the Euclidean result in each chart and take a subsequence common to this finite collection. Since the cutoffs sum to one and their volume densities are uniformly comparable, the resulting sequence is Cauchy in $L^2(M)$. This proves the compactness statement without imposing a bound on the values of $f_j$.
\end{proof}

\begin{corollary}[A positive Poincare constant on a connected closed manifold]\label{cor:closed-manifold-poincare}
If $M$ is additionally nonempty, connected and positive-dimensional, then
\[
 \lambda_1=\inf_{\substack{v\in H^1(M),\ \int_Mv=0\\v\ne0}}
      \frac{\int_M|\nabla v|^2}{\int_Mv^2}>0.
\]
Consequently $\lambda_1\int_M|f-\bar f|^2\leq\int_M|\nabla f|^2$ for every $f\in H^1(M)$, where $\bar f$ is its average.
\end{corollary}
\begin{proof}
If the infimum were zero, choose mean-zero $v_j$ of $L^2$ norm one with $\|\nabla v_j\|_2\to0$. Theorem \ref{thm:sobolev-closed-l2-compactness} supplies a strongly $L^2$ convergent subsequence. Distributional testing in each chart shows that the limit has zero weak gradient. Lemma \ref{lem:weak-gradient-zero} makes it constant locally, and connectedness propagates that constant over the manifold. Its mean is zero, so it is zero, contradicting the limiting $L^2$ norm one. The inequality is then the Rayleigh definition applied to $f-\bar f$. On a disconnected manifold, separate constants on different components show why a single global mean-zero condition would not suffice.
\end{proof}
